\documentclass[11pt,a4paper]{article}
\usepackage[margin=0.9in]{geometry}
\usepackage{amsmath,amssymb,amsthm}
\usepackage{graphicx}
\usepackage{xcolor}
\usepackage{mdframed}
\usepackage{enumitem}
\usepackage{booktabs}
\usepackage{titlesec}
\usepackage[colorlinks=true,linkcolor=blue!50!black,urlcolor=blue!50!black]{hyperref}
\usepackage{fancyhdr}
\usepackage{tikz}
\usetikzlibrary{arrows.meta,positioning}

\definecolor{keyblue}{RGB}{22,78,128}
\definecolor{warmred}{RGB}{160,40,40}
\definecolor{boxbg}{RGB}{242,246,250}
\definecolor{trapbg}{RGB}{253,244,244}
\definecolor{exbg}{RGB}{244,250,244}

\titleformat{\section}{\Large\bfseries\color{keyblue}}{\thesection}{0.7em}{}
\titleformat{\subsection}{\large\bfseries\color{keyblue!85!black}}{\thesubsection}{0.6em}{}

\pagestyle{fancy}
\fancyhf{}
\fancyhead[L]{\small\color{keyblue}EE341 Prereq Revamp — Signals, Systems \& Fourier}
\fancyfoot[C]{\small\thepage}
\renewcommand{\headrulewidth}{0.4pt}

\newmdenv[backgroundcolor=boxbg,linecolor=keyblue,linewidth=1pt,
  leftmargin=0pt,rightmargin=0pt,innerleftmargin=8pt,innerrightmargin=8pt,
  innertopmargin=6pt,innerbottommargin=6pt,skipabove=8pt,skipbelow=8pt]{keybox}
\newmdenv[backgroundcolor=trapbg,linecolor=warmred,linewidth=1pt,
  leftmargin=0pt,rightmargin=0pt,innerleftmargin=8pt,innerrightmargin=8pt,
  innertopmargin=6pt,innerbottommargin=6pt,skipabove=8pt,skipbelow=8pt]{trapbox}
\newmdenv[backgroundcolor=exbg,linecolor=green!35!black,linewidth=1pt,
  leftmargin=0pt,rightmargin=0pt,innerleftmargin=8pt,innerrightmargin=8pt,
  innertopmargin=6pt,innerbottommargin=6pt,skipabove=8pt,skipbelow=8pt]{exbox}

\newcommand{\FT}{\;\stackrel{\mathcal{F}}{\longleftrightarrow}\;}
\newcommand{\sinc}{\operatorname{sinc}}
\newcommand{\rect}{\operatorname{rect}}
\newcommand{\tri}{\operatorname{tri}}
\newcommand{\sgn}{\operatorname{sgn}}
\newcommand{\dd}{\,\mathrm{d}}
\newcommand{\ii}{\mathrm{j}}

\newenvironment{worked}[1]{\par\vspace{4pt}\noindent\textbf{\color{warmred}Worked problem #1.}\ \ignorespaces}{\par\vspace{4pt}}
\newcommand{\soln}{\par\vspace{2pt}\noindent\textit{Solution.}\ \ignorespaces}

\title{\vspace{-1.5cm}\textbf{\color{keyblue}Signals, Systems \& the Fourier Transform}\\[4pt]
\large A prerequisite revamp for EE341 (Communication Systems--I), IIT Bombay}
\author{}
\date{}

\begin{document}
\maketitle
\vspace{-1.2cm}

\begin{keybox}
\textbf{How to use this.} You said your fundamentals are shaky, so this builds each idea up rather than
just listing formulas. But your first quiz is \textbf{17 August}, so it is ordered by \emph{what EE341 actually
uses}, not by what a full signals course would cover.

\medskip
\noindent\textbf{If you have three days:}
\begin{itemize}[nosep,leftmargin=*]
\item \textbf{Day 1} --- Part A (\S1--\S3) and Part B (\S4--\S6). This is the load-bearing material. Do every worked problem \emph{yourself} first.
\item \textbf{Day 2} --- Part C (\S7--\S9). Sampling is standard; \S9 (bandpass / lowpass-equivalent) is where your lecture notes already are, so it matters most.
\item \textbf{Day 3} --- Part D (\S10) skim only, then the problem set in \S11 under time pressure.
\end{itemize}

\medskip
\noindent\textbf{Notation warning.} EE341 (and Madhow's book) uses \textbf{ordinary frequency $f$ in Hz}, not
$\omega$ in rad/s. Your notes write $X(f)$, $\cos(2\pi f_c t)$, never $\omega$. Everything here follows that
convention. It removes all the annoying $1/2\pi$ factors --- the transform becomes perfectly symmetric.
\end{keybox}

\tableofcontents
\newpage

%==============================================================
\section*{\color{keyblue}PART A --- Signals and Systems}
\addcontentsline{toc}{section}{PART A --- Signals and Systems}
%==============================================================

\section{What a signal is, and the four classifications that matter}

A \textbf{signal} is a function carrying information. In this course it is almost always $x(t)$, a real-valued
function of continuous time. Four classifications come up constantly in communications, and each one is really
a question about \emph{which mathematical tool applies}.

\subsection{Energy vs.\ power signals}

Define
\[
E_x = \int_{-\infty}^{\infty}|x(t)|^2\dd t,
\qquad
P_x = \lim_{T\to\infty}\frac{1}{T}\int_{-T/2}^{T/2}|x(t)|^2\dd t .
\]

\begin{itemize}[leftmargin=*]
\item If $0<E_x<\infty$, $x$ is an \textbf{energy signal}. Then $P_x=0$ automatically. Pulses, decaying
exponentials, anything that dies out.
\item If $E_x=\infty$ but $0<P_x<\infty$, $x$ is a \textbf{power signal}. Sinusoids, periodic signals, random
noise, and \emph{every real communication waveform} (a transmitter runs forever).
\end{itemize}

\begin{keybox}
\textbf{Why you should care.} This determines which spectral object exists.
An energy signal has a Fourier transform $X(f)$ and an \emph{energy} spectral density $|X(f)|^2$.
A power signal generally does \emph{not} have a well-behaved $X(f)$ --- it has a \emph{power} spectral
density instead. This is exactly why Part D (random processes) exists: modulated signals and noise are power
signals, so ``take the Fourier transform of the noise'' is meaningless, and you must talk about PSD.
\end{keybox}

\begin{worked}{1.1}
Your notes (page 4) compute the power of the DSB-SC signal $x_c(t)=A\,x(t)\cos(2\pi f_c t)$. Reproduce it.
\end{worked}
\soln
\[
P_{x_c}=\lim_{T\to\infty}\frac1T\int_{-T/2}^{T/2}A^2x^2(t)\cos^2(2\pi f_c t)\dd t
=\lim_{T\to\infty}\frac1T\int_{-T/2}^{T/2}A^2x^2(t)\,\frac{1+\cos(4\pi f_c t)}{2}\dd t .
\]
Split it. The first piece gives $\tfrac{A^2}{2}P_x$. The second piece,
$\tfrac{A^2}{2}\langle x^2(t)\cos(4\pi f_ct)\rangle$, vanishes: $x^2(t)$ is lowpass with bandwidth $\le 2W$,
while $\cos(4\pi f_c t)$ oscillates at $2f_c\gg 2W$, so the product averages to zero. Hence
\[
\boxed{P_{x_c}=\frac{A^2}{2}P_x}.
\]
For DSB-AM, $x_c'(t)=A(1+\mu x(t))\cos(2\pi f_ct)$, the same argument on $(1+\mu x)^2 = 1+2\mu x+\mu^2x^2$
with $\langle x\rangle=0$ gives
\[
P_{x_c'}=\frac{A^2}{2}\big(1+\mu^2 P_x\big)=\underbrace{\frac{A^2}{2}}_{\text{carrier, useless}}+\underbrace{\frac{A^2\mu^2}{2}P_x}_{\text{message}} .
\]
The \textbf{power efficiency} $\eta=\mu^2P_x/(1+\mu^2P_x)$ is why DSB-AM is wasteful --- a favourite quiz question.

\subsection{Periodic vs.\ aperiodic}
$x(t)$ is periodic if $x(t+T_0)=x(t)$ for some $T_0>0$; the smallest such $T_0$ is the fundamental period,
$f_0=1/T_0$. Periodic $\Rightarrow$ Fourier \emph{series} (discrete spectrum, spectral lines). Aperiodic
$\Rightarrow$ Fourier \emph{transform} (continuous spectrum). \S4 shows these are the same theorem.

\subsection{Real vs.\ complex}
Physical signals are real. This is not a triviality --- it forces \textbf{Hermitian symmetry} on the spectrum
($\S5.4$), which is the single most-used structural fact in this course. Complex signals appear as
\emph{bookkeeping devices}: the analytic signal $x_+(t)$ and the complex envelope $x_\ell(t)$ in \S9 are complex
by construction, and that is the whole point.

\subsection{Causal vs.\ non-causal; lowpass vs.\ bandpass}
Causal: $x(t)=0$ for $t<0$. Matters for realizability of filters.

Lowpass vs.\ bandpass is the classification your course opens with. From your notes, page 1:
\begin{keybox}
$x(t)$ is a \textbf{lowpass signal of bandwidth $W$} if $X(f)=0$ for all $|f|>W$.\\
$x(t)$ is a \textbf{bandpass signal} with centre $f_c$ and bandwidth $2W$ if $X(f)=0$ except for
$f_c-W\le|f|\le f_c+W$.
\end{keybox}
\textbf{Careful with the word ``bandwidth''.} For a real lowpass signal occupying $[-W,W]$, the bandwidth is
$W$ (the one-sided extent), not $2W$. For the bandpass signal occupying $[f_c-W,f_c+W]$ and its mirror, the
bandwidth is $2W$. Your notes cross out a ``$2$'' on page 1 for exactly this reason. Sloppiness here loses
marks; state which convention you are using.

\section{Elementary signals you must know cold}

\begin{center}
\renewcommand{\arraystretch}{1.5}
\begin{tabular}{@{}lll@{}}
\toprule
\textbf{Name} & \textbf{Definition} & \textbf{Note} \\
\midrule
Unit rectangle & $\rect(t)=1$ for $|t|<\tfrac12$, else $0$ & width 1, centred at 0 \\
Unit triangle & $\tri(t)=1-|t|$ for $|t|<1$, else $0$ & $=\rect*\rect$ \\
Sinc & $\sinc(t)=\dfrac{\sin(\pi t)}{\pi t}$, $\sinc(0)=1$ & zeros at all nonzero integers \\
Unit step & $u(t)=1$ for $t>0$, $0$ for $t<0$, $\tfrac12$ at $0$ & \\
Signum & $\sgn(t)=2u(t)-1$ & \\
Impulse & $\delta(t)$, see below & not a function \\
\bottomrule
\end{tabular}
\end{center}

\begin{trapbox}
\textbf{Trap.} There are two sinc conventions. Communications uses the \textbf{normalised} one,
$\sinc(t)=\sin(\pi t)/(\pi t)$, so that $\rect(t)\FT\sinc(f)$ cleanly. Some maths and DSP texts use
$\sin(t)/t$. Always check. Every formula in this document uses the normalised one.
\end{trapbox}

\subsection{The delta function, properly}
$\delta(t)$ is not a function; it is defined \emph{entirely} by how it behaves inside an integral:
\[
\int_{-\infty}^{\infty}\varphi(t)\,\delta(t-t_0)\dd t=\varphi(t_0)\qquad\text{(sifting property)}
\]
for every continuous $\varphi$. Everything else follows from this. Two consequences you will use daily:
\begin{align}
x(t)*\delta(t-t_0)&=x(t-t_0) &&\text{(convolving with a shifted impulse = shifting)}\\
x(t)\,\delta(t-t_0)&=x(t_0)\,\delta(t-t_0) &&\text{(multiplying by an impulse = sampling)}
\end{align}
Also the scaling rule $\delta(at)=\frac{1}{|a|}\delta(t)$, which trips people up constantly.

\begin{trapbox}
\textbf{Trap.} $\delta(t)$ has \emph{unit area}, not unit height. So in a spectrum plot an impulse is drawn as
an arrow whose labelled height is its \emph{weight} (area). When your notes draw the carrier line in the
DSB-AM spectrum as a tall arrow at $f_c$, that arrow has weight $A/2$, i.e.\ the term
$\frac{A}{2}\delta(f-f_c)$ --- it is not ``infinitely large signal''.
\end{trapbox}

\section{LTI systems and convolution}

A system $\mathcal{H}$ maps input $x(t)$ to output $y(t)$. It is
\begin{itemize}[nosep,leftmargin=*]
\item \textbf{Linear} if $\mathcal{H}\{a x_1+b x_2\}=a\mathcal{H}\{x_1\}+b\mathcal{H}\{x_2\}$;
\item \textbf{Time-invariant} if $\mathcal{H}\{x(t-t_0)\}=y(t-t_0)$.
\end{itemize}
Both together: \textbf{LTI}.

\subsection{Why LTI is the whole game}
Here is the argument you should be able to reproduce, because it explains why Fourier analysis exists at all.

Write any input using the sifting property as a continuum of shifted impulses:
\[
x(t)=\int_{-\infty}^{\infty}x(\tau)\,\delta(t-\tau)\dd\tau .
\]
Let $h(t)=\mathcal{H}\{\delta(t)\}$ be the \textbf{impulse response}. By time-invariance,
$\mathcal{H}\{\delta(t-\tau)\}=h(t-\tau)$. By linearity (extended to integrals), push $\mathcal H$ through:
\[
\boxed{\;y(t)=\int_{-\infty}^{\infty}x(\tau)\,h(t-\tau)\dd\tau \;=\; (x*h)(t)\;}
\]
So \textbf{an LTI system is completely characterised by one function, $h(t)$}. That is an enormous reduction,
and it is why we can talk about ``the channel'' or ``the filter'' as a single object.

\subsection{Properties of convolution}
\begin{align*}
&\text{Commutative: } x*h=h*x
&&\text{Associative: } (x*h_1)*h_2=x*(h_1*h_2)\\
&\text{Distributive: } x*(h_1+h_2)=x*h_1+x*h_2
&&\text{Shift: } x(t)*h(t-t_0)=y(t-t_0)
\end{align*}
Associativity is why cascaded filters multiply their frequency responses; distributivity is why parallel
branches add.

\textbf{Widths add.} If $x$ is supported on $[a_1,b_1]$ and $h$ on $[a_2,b_2]$, then $x*h$ is supported on
$[a_1+a_2,\,b_1+b_2]$. Use this as an instant sanity check on any convolution you compute.

\begin{worked}{3.1}
Compute $\rect(t)*\rect(t)$.
\end{worked}
\soln
Both have width 1, so the answer has width 2, supported on $[-1,1]$ --- already a check.
\[
(\rect*\rect)(t)=\int_{-\infty}^{\infty}\rect(\tau)\rect(t-\tau)\dd\tau
=\int_{-1/2}^{1/2}\rect(t-\tau)\dd\tau .
\]
The integrand is 1 when $|t-\tau|<\tfrac12$, i.e.\ $\tau\in(t-\tfrac12,\,t+\tfrac12)$. So the integral is the
length of the overlap between $(-\tfrac12,\tfrac12)$ and $(t-\tfrac12,t+\tfrac12)$, which is $1-|t|$ for
$|t|\le1$ and 0 otherwise. Therefore
\[
\rect(t)*\rect(t)=\tri(t).
\]
This one result will save you repeatedly: it is why $\sinc^2$ pairs with the triangle, and it is the reason a
raised-cosine-type pulse shows up later in the ISI chapter.

\subsection{Causality and stability from $h(t)$}
\begin{itemize}[nosep,leftmargin=*]
\item \textbf{Causal} $\iff h(t)=0$ for $t<0$.
\item \textbf{BIBO stable} $\iff \int|h(t)|\dd t<\infty$.
\end{itemize}
The ideal lowpass filter has $h(t)=2W\sinc(2Wt)$, which is neither causal nor absolutely integrable. That is
why ``ideal filter'' is a mathematical fiction we use for analysis and approximate in practice --- worth one
sentence in any exam answer that invokes one.

\subsection{Eigenfunctions: the bridge to Fourier}
Feed a complex exponential $x(t)=e^{\ii 2\pi f_0 t}$ into an LTI system:
\[
y(t)=\int h(\tau)e^{\ii2\pi f_0 (t-\tau)}\dd\tau
= e^{\ii2\pi f_0 t}\underbrace{\int h(\tau)e^{-\ii2\pi f_0\tau}\dd\tau}_{\textstyle H(f_0)} = H(f_0)\,e^{\ii2\pi f_0t}.
\]
\begin{keybox}
\textbf{The complex exponential comes out unchanged except for a complex scale factor.} It is an
\emph{eigenfunction} of every LTI system, with eigenvalue $H(f_0)$ --- and that eigenvalue is precisely the
Fourier transform of $h$.

This is \emph{the} reason we decompose signals into complex exponentials rather than, say, polynomials:
because it turns the hard operation (convolution) into the easy one (multiplication). Everything in Part B is
machinery for doing that decomposition.
\end{keybox}

\begin{trapbox}
\textbf{Trap.} $\cos(2\pi f_0t)$ is \emph{not} an eigenfunction --- an LTI system can change its phase, so the
output is $|H(f_0)|\cos(2\pi f_0 t+\angle H(f_0))$, not a scalar multiple of the input. Only the complex
exponential is a true eigenfunction. (For a real system the cosine result follows by writing it as a sum of
two exponentials and using $H(-f)=H^*(f)$.)
\end{trapbox}
%==============================================================
\section*{\color{keyblue}PART B --- The Fourier Transform}
\addcontentsline{toc}{section}{PART B --- The Fourier Transform}
%==============================================================

\section{From Fourier series to Fourier transform}

\subsection{The series}
Any reasonable periodic signal with period $T_0$ ($f_0=1/T_0$) can be written as
\[
x(t)=\sum_{n=-\infty}^{\infty}c_n e^{\ii2\pi n f_0 t},
\qquad
c_n=\frac{1}{T_0}\int_{T_0}x(t)e^{-\ii2\pi n f_0 t}\dd t .
\]
The $c_n$ are complex numbers: $|c_n|$ is the amplitude of the $n$-th harmonic, $\angle c_n$ its phase. The
spectrum is a set of \textbf{discrete lines} at $f=nf_0$.

Why does the formula for $c_n$ work? Because the exponentials are orthogonal over a period:
$\frac{1}{T_0}\int_{T_0}e^{\ii2\pi(n-m)f_0t}\dd t=\delta_{nm}$. Multiplying the series by $e^{-\ii2\pi mf_0t}$
and averaging kills every term except $n=m$. That is the entire trick, and the same trick reappears in matched
filtering and in orthogonal signalling later in the course.

\begin{worked}{4.1}
Find the Fourier series of the square wave $c(t)$ your notes use for the ring modulator / switching
demodulator: $c(t)=+1$ for the first half of each period, $-1$ for the second, period $T_c=1/f_c$.
\end{worked}
\soln
$c(t)$ is odd (with the origin chosen at a rising edge midpoint) and has zero mean, so only odd harmonics
survive:
\[
c(t)=\frac{4}{\pi}\left[\cos(2\pi f_ct)-\frac13\cos(6\pi f_ct)+\frac15\cos(10\pi f_ct)-\cdots\right].
\]
That matches your notes exactly ($\frac4\pi\cos\omega_ct-\frac{4}{3\pi}\cos3\omega_ct+\frac{4}{5\pi}\cos5\omega_ct$).
The \emph{unipolar} switching function (0/1 instead of $\mp1$) that appears on the demodulation page is just
$\frac{1+c(t)}{2}$:
\[
s(t)=\frac12+\frac{2}{\pi}\cos(2\pi f_ct)-\frac{2}{3\pi}\cos(6\pi f_ct)+\cdots
\]
--- again exactly your notes. \textbf{Why this matters:} a switching modulator multiplies $x(t)$ by $s(t)$, so
the output contains $x(t)$ replicated around DC, $f_c$, $3f_c$, $5f_c,\dots$ A bandpass filter at $f_c$ with
bandwidth $2W$ picks off the term $\frac{2}{\pi}x(t)\cos(2\pi f_ct)$, which is DSB-SC. Cheap hardware, exact
result.

\subsection{Letting the period go to infinity}
Take an aperiodic $x(t)$, and build a periodic version $\tilde x_{T_0}$ by repeating it every $T_0$. Then
\[
c_n=\frac{1}{T_0}\int_{-T_0/2}^{T_0/2}x(t)e^{-\ii2\pi nf_0t}\dd t=\frac{1}{T_0}X(nf_0),
\qquad X(f):=\int_{-\infty}^{\infty}x(t)e^{-\ii2\pi ft}\dd t .
\]
The lines sit at spacing $f_0=1/T_0$ and have height $X(nf_0)/T_0$. As $T_0\to\infty$ the lines merge into a
continuum, the sum becomes an integral, and we get the transform pair.

\begin{keybox}
\textbf{The Fourier transform pair (Hz convention).}
\[
X(f)=\int_{-\infty}^{\infty}x(t)\,e^{-\ii2\pi f t}\dd t,
\qquad
x(t)=\int_{-\infty}^{\infty}X(f)\,e^{+\ii2\pi f t}\dd f .
\]
Note the beautiful symmetry: no $1/2\pi$ anywhere, and forward/inverse differ only in the sign of the exponent.
This is the payoff for using $f$ instead of $\omega$.
\end{keybox}

\begin{trapbox}
\textbf{Converting from an $\omega$-based textbook.} If a book gives $X(\omega)$, then $X(f)=X(\omega)\big|_{\omega=2\pi f}$
--- the \emph{function values} are the same, only the argument is relabelled. But a formula like
$x(t)=\frac{1}{2\pi}\int X(\omega)e^{\ii\omega t}\dd\omega$ carries a $1/2\pi$ that disappears in the $f$
convention. Never mix the two mid-problem.
\end{trapbox}

\subsection{What $X(f)$ actually means}
$X(f)$ is complex. Write $X(f)=|X(f)|e^{\ii\theta(f)}$.
\begin{itemize}[nosep,leftmargin=*]
\item $|X(f)|$ --- \textbf{magnitude spectrum}: how much of frequency $f$ is present.
\item $\theta(f)$ --- \textbf{phase spectrum}: \emph{where} those components sit in time. Phase carries most
of the perceptual information in a signal; discarding it destroys the waveform even though $|X|$ is intact.
\item Units: if $x$ is in volts, $X(f)$ is in volts/Hz. It is a \emph{density}, which is why impulses appear
whenever a signal has finite power concentrated at a single frequency.
\end{itemize}

\section{Properties of the Fourier transform}

These are not trivia to memorise --- in this course you will almost never compute an FT integral. You will
recognise a signal as a combination of known pairs and apply properties. Learn them as \emph{moves}.

\subsection{Linearity}
$a x_1(t)+b x_2(t)\FT aX_1(f)+bX_2(f)$. Obvious, but it is what lets you decompose DSB-AM into
carrier + sidebands and treat each separately.

\subsection{Time shift and frequency shift (the two you will use most)}
\begin{align}
x(t-t_0)&\FT X(f)\,e^{-\ii2\pi ft_0} &&\text{delay $\Rightarrow$ linear phase, magnitude unchanged}\\
x(t)\,e^{\ii2\pi f_ct}&\FT X(f-f_c) &&\text{multiply by exponential $\Rightarrow$ shift spectrum}
\end{align}
\textbf{Proof of the second} (do this once by hand, it takes two lines):
$\int x(t)e^{\ii2\pi f_ct}e^{-\ii2\pi ft}\dd t=\int x(t)e^{-\ii2\pi(f-f_c)t}\dd t=X(f-f_c)$.

\begin{keybox}
\textbf{The modulation property --- the single most important formula in this course.}
\[
x(t)\cos(2\pi f_ct)\;=\;\frac{x(t)e^{\ii2\pi f_ct}+x(t)e^{-\ii2\pi f_ct}}{2}
\FT
\frac{1}{2}X(f-f_c)+\frac{1}{2}X(f+f_c).
\]
Multiplying by a carrier \textbf{takes the spectrum, halves it, and puts a copy at $+f_c$ and a copy at
$-f_c$}. That is literally the definition of modulation as your notes state it: ``process of shifting
frequency''. Everything in the AM chapter is this one line applied repeatedly.
\end{keybox}

Your notes write $X_c(f)=\frac{A}{2}X(f-f_c)+\frac{A}{2}X(f+f_c)$ for $x_c(t)=Ax(t)\cos(2\pi f_ct)$, and then
say ``$=\frac{A}{2}X(f-f_c)$ for $f>0$, and $X_c(f)$ has Hermitian symmetry''. That shorthand is legitimate
\emph{only} because $X(f-f_c)$ and $X(f+f_c)$ do not overlap, which needs $f_c>W$. State that condition.

\subsection{Scaling}
\[
x(at)\FT\frac{1}{|a|}X\!\left(\frac fa\right).
\]
\textbf{Compress in time $\Rightarrow$ stretch in frequency.} This is the time--bandwidth trade-off: a short
pulse necessarily occupies a wide band. It is the reason high data rates need wide bandwidth, which is the
theme of the second half of EE341.

\subsection{Duality}
If $x(t)\FT X(f)$, then $X(t)\FT x(-f)$. This gives you a free pair for every pair you know:
from $\rect(t)\FT\sinc(f)$ you immediately get $\sinc(t)\FT\rect(f)$ (using evenness). Use it; it halves
what you need to memorise.

\subsection{Conjugate symmetry --- Hermitian symmetry}
\begin{keybox}
If $x(t)$ is \textbf{real}, then
\[
\boxed{X(-f)=X^*(f)}
\]
equivalently $|X(f)|$ is \textbf{even} and $\angle X(f)$ is \textbf{odd}.
\end{keybox}
\textbf{Proof.} $X(-f)=\int x(t)e^{\ii2\pi ft}\dd t = \left[\int x^*(t)e^{-\ii2\pi ft}\dd t\right]^*=X^*(f)$
since $x=x^*$.

\textbf{Why this dominates the course.} It says the negative-frequency half of the spectrum of a real signal
carries \emph{no new information}. Hence:
\begin{itemize}[nosep,leftmargin=*]
\item You may specify a real bandpass signal by its positive-frequency content alone --- this is what your
notes mean by ``$X_c(f)=A X(f-f_c)$ for $f>0$, and $X_c(f)$ has Hermitian symmetry''.
\item It motivates the \textbf{analytic signal} $x_+(t)$ (\S9), which simply \emph{deletes} the redundant half.
\item It is why DSB transmission is spectrally wasteful (upper and lower sidebands are mirror images), which
motivates SSB and VSB.
\end{itemize}
Two corollaries worth having reflexively: $x$ real and even $\Rightarrow$ $X$ real and even;
$x$ real and odd $\Rightarrow$ $X$ purely imaginary and odd.

\subsection{Convolution and multiplication}
\begin{keybox}
\[
x(t)*h(t)\FT X(f)H(f)
\qquad\qquad
x(t)\,h(t)\FT X(f)*H(f)
\]
\end{keybox}
The first is the reason Fourier analysis is used at all (\S3.5). The second is its dual and is how you handle
any product of signals --- e.g.\ the square-law modulator, where you square $x(t)+\cos2\pi f_ct$ and want to
know where the resulting spectral pieces land.

\textbf{Proof of the first} (know it):
\[
\int\!\!\left[\int x(\tau)h(t-\tau)\dd\tau\right]e^{-\ii2\pi ft}\dd t
=\int x(\tau)\!\left[\int h(t-\tau)e^{-\ii2\pi ft}\dd t\right]\!\dd\tau
=\int x(\tau)H(f)e^{-\ii2\pi f\tau}\dd\tau=X(f)H(f),
\]
where the inner bracket used the time-shift property.

\subsection{Differentiation and integration}
\[
\frac{\dd x}{\dd t}\FT \ii2\pi f\,X(f),
\qquad
\int_{-\infty}^{t}x(\tau)\dd\tau\FT\frac{X(f)}{\ii2\pi f}+\frac{X(0)}{2}\delta(f).
\]
The $\delta(f)$ term (DC offset) is forgotten by almost everyone. It is needed exactly when $X(0)\ne0$.

\subsection{Parseval / Rayleigh}
\[
\int_{-\infty}^{\infty}|x(t)|^2\dd t=\int_{-\infty}^{\infty}|X(f)|^2\dd f
\qquad\text{(energy signals)}
\]
More generally $\int x(t)y^*(t)\dd t=\int X(f)Y^*(f)\dd f$. The quantity $|X(f)|^2$ is the
\textbf{energy spectral density} in J/Hz. For \emph{power} signals the analogue is the power spectral density
of \S10 --- do not apply this Parseval to a modulated carrier.

\subsection{Moment / area rules (fast sanity checks)}
\[
X(0)=\int x(t)\dd t,
\qquad
x(0)=\int X(f)\dd f .
\]
Use these constantly. If you compute a transform and $X(0)$ does not equal the area under $x(t)$, you have
made an algebra error. Thirty seconds, catches most mistakes.

\section{Standard transform pairs}

\begin{center}
\renewcommand{\arraystretch}{1.7}
\begin{tabular}{@{}p{6.6cm}p{6.6cm}@{}}
\toprule
\textbf{$x(t)$} & \textbf{$X(f)$}\\
\midrule
$\delta(t)$ & $1$\\
$1$ & $\delta(f)$\\
$\delta(t-t_0)$ & $e^{-\ii2\pi ft_0}$\\
$e^{\ii2\pi f_ct}$ & $\delta(f-f_c)$\\
$\cos(2\pi f_ct)$ & $\tfrac12\delta(f-f_c)+\tfrac12\delta(f+f_c)$\\
$\sin(2\pi f_ct)$ & $\tfrac{1}{2\ii}\delta(f-f_c)-\tfrac{1}{2\ii}\delta(f+f_c)$\\
$\rect(t/T)$ & $T\sinc(fT)$\\
$2W\sinc(2Wt)$ & $\rect\!\big(f/(2W)\big)$\\
$\tri(t/T)$ & $T\sinc^2(fT)$\\
$e^{-at}u(t)$, $a>0$ & $\dfrac{1}{a+\ii2\pi f}$\\
$e^{-a|t|}$, $a>0$ & $\dfrac{2a}{a^2+(2\pi f)^2}$\\
$e^{-\pi t^2}$ & $e^{-\pi f^2}$ \quad (its own transform)\\
$u(t)$ & $\dfrac{1}{2}\delta(f)+\dfrac{1}{\ii2\pi f}$\\
$\sgn(t)$ & $\dfrac{1}{\ii\pi f}$\\
$\dfrac{1}{\pi t}$ & $-\ii\,\sgn(f)$ \quad \textbf{(Hilbert transform kernel)}\\
$\displaystyle\sum_{n}\delta(t-nT_s)$ & $\dfrac{1}{T_s}\displaystyle\sum_{k}\delta\!\left(f-\tfrac{k}{T_s}\right)$ \quad \textbf{(impulse train)}\\
\bottomrule
\end{tabular}
\end{center}

\begin{keybox}
\textbf{The three lines to memorise above all others}, because EE341 uses them on nearly every page:
\begin{enumerate}[nosep,leftmargin=*]
\item $\cos(2\pi f_ct)\FT\tfrac12[\delta(f-f_c)+\delta(f+f_c)]$ --- modulation.
\item $\dfrac{1}{\pi t}\FT-\ii\sgn(f)$ --- Hilbert transform, hence analytic signal and SSB.
\item $\sum_n\delta(t-nT_s)\FT \dfrac{1}{T_s}\sum_k\delta(f-k/T_s)$ --- sampling theorem.
\end{enumerate}
\end{keybox}

\begin{worked}{6.1}
Derive $u(t)\FT\frac12\delta(f)+\frac{1}{\ii2\pi f}$, and use it to get your notes' expression
$U(f)=\frac12(1+\sgn(f))$ in the frequency variable.
\end{worked}
\soln
Write $u(t)=\frac12+\frac12\sgn(t)$. The constant $\frac12$ transforms to $\frac12\delta(f)$.
For $\sgn(t)$: it is not integrable, so take the limit of $e^{-a|t|}\sgn(t)$ as $a\to0^+$:
\[
\int_0^\infty e^{-at}e^{-\ii2\pi ft}\dd t-\int_{-\infty}^0 e^{at}e^{-\ii2\pi ft}\dd t
=\frac{1}{a+\ii2\pi f}-\frac{1}{a-\ii2\pi f}\xrightarrow{a\to0^+}\frac{-\ii4\pi f}{(2\pi f)^2}=\frac{1}{\ii\pi f}.
\]
Hence $u(t)\FT\frac12\delta(f)+\frac{1}{\ii2\pi f}$.

Now the notes' $U(f)$ is a \emph{different object}: it is the unit step \emph{as a function of frequency},
used as a mask to kill negative frequencies. There, $U(f)=\frac12(1+\sgn(f))$ is simply the algebraic identity
$u=\frac12+\frac12\sgn$ written with argument $f$: $U(f)=0$ for $f<0$, $\frac12$ at $f=0$, $1$ for $f>0$ ---
precisely the box on your SNR page. Keep the two straight: one is a transform, the other is a masking
function.

\begin{worked}{6.2}
A pulse $p(t)=\rect(t/T)$ is transmitted. Find its bandwidth by the first-null definition, and comment on the
time--bandwidth product.
\end{worked}
\soln
$P(f)=T\sinc(fT)$, first null at $f=1/T$. So $B_{\text{null}}=1/T$ and $B\cdot T=1$. Halving the pulse
duration doubles the bandwidth --- the scaling property in action. This constant is the seed of the Nyquist
criterion in the ISI chapter: you cannot signal faster than roughly $2B$ symbols/second in bandwidth $B$.
%==============================================================
\section*{\color{keyblue}PART C --- Filtering, Sampling, and Bandpass Signals}
\addcontentsline{toc}{section}{PART C --- Filtering, Sampling, Bandpass Signals}
%==============================================================

\section{Filtering in the frequency domain}

For an LTI system, $Y(f)=H(f)X(f)$ where $H(f)$ is the \textbf{frequency response} (the FT of $h(t)$).
Magnitude $|H(f)|$ scales each frequency; phase $\angle H(f)$ delays it.

\subsection{Distortionless transmission}
A channel is distortionless if $y(t)=K\,x(t-t_d)$, i.e.
\[
H(f)=K e^{-\ii2\pi f t_d}
\quad\Longleftrightarrow\quad
|H(f)|=K \text{ (flat)}, \qquad \angle H(f)=-2\pi f t_d \text{ (linear in $f$)}.
\]
\begin{keybox}
\textbf{Flat magnitude and \emph{linear} phase} over the signal's band. Nonflat magnitude gives amplitude
distortion; nonlinear phase gives \textbf{delay distortion} --- different frequency components arrive at
different times, smearing pulses. This is the same phenomenon that becomes intersymbol interference later.
The \emph{group delay} $\tau_g(f)=-\frac{1}{2\pi}\frac{\dd\angle H(f)}{\dd f}$ is constant exactly when phase
is linear.
\end{keybox}

\subsection{Ideal filters}
\begin{itemize}[leftmargin=*]
\item \textbf{Ideal LPF}, cutoff $W$: $H(f)=\rect\!\big(\tfrac{f}{2W}\big)$, so $h(t)=2W\sinc(2Wt)$.
\item \textbf{Ideal BPF}, centre $f_c$, bandwidth $2W$: $H(f)=\rect\!\big(\tfrac{f-f_c}{2W}\big)+\rect\!\big(\tfrac{f+f_c}{2W}\big)$,
so $h(t)=4W\sinc(2Wt)\cos(2\pi f_ct)$ --- a lowpass sinc \emph{modulated} onto the carrier. Note how the
modulation property gives you the bandpass impulse response for free.
\end{itemize}
Both are non-causal (the sinc extends to $t<0$) and unstable in the BIBO sense. Real filters (Butterworth,
Chebyshev) approximate them with finite roll-off, and a real receiver's ``high-$Q$ tunable bandpass filter''
--- the thing your superheterodyne page says is hard to build --- is exactly this approximation problem.

\begin{worked}{7.1}
In the coherent demodulator on your notes page 7: $x_c'(t)=Ax(t)\cos(2\pi f_ct)$ is multiplied by
$\cos(2\pi f_ct+\phi)$ and lowpass filtered. Find the output.
\end{worked}
\soln
\[
y(t)=Ax(t)\cos(2\pi f_ct)\cos(2\pi f_ct+\phi)=\frac{Ax(t)}{2}\big[\cos\phi+\cos(4\pi f_ct+\phi)\big].
\]
The LPF (cutoff $W<2f_c-W$) removes the $2f_c$ term, leaving
\[
\boxed{y_{\text{out}}(t)=\frac{A\cos\phi}{2}\,x(t)}.
\]
\begin{keybox}
\textbf{The punchline of coherent detection.} The recovered signal is scaled by $\cos\phi$. If the local
oscillator is $90^\circ$ out of phase ($\phi=\pi/2$) the output is \emph{zero} --- total loss of signal. This
is the \textbf{quadrature null effect}, and it is why DSB-SC needs carrier recovery (a PLL), and why DSB-AM
with envelope detection --- which needs no phase reference at all --- survived commercially despite being
power-inefficient.
\end{keybox}
If instead there is a frequency error $\Delta f$ (LO at $f_c+\Delta f$), you get
$\frac{A}{2}x(t)\cos(2\pi\Delta f\,t)$: a slow beating in and out. Your notes' page 11 does exactly this
calculation in complex-envelope form.

\section{Sampling}

\subsection{Impulse sampling}
Sample $x(t)$ every $T_s$ seconds. Model it as multiplication by an impulse train:
\[
x_s(t)=x(t)\sum_{n}\delta(t-nT_s)=\sum_n x(nT_s)\delta(t-nT_s).
\]
Multiplication in time $\Rightarrow$ convolution in frequency, and the impulse train transforms to another
impulse train with spacing $f_s=1/T_s$:
\[
X_s(f)=X(f)*\frac{1}{T_s}\sum_k\delta\!\left(f-kf_s\right)
=\boxed{f_s\sum_{k=-\infty}^{\infty}X(f-kf_s)}
\]

\begin{keybox}
\textbf{Sampling replicates the spectrum} at every multiple of $f_s$, scaled by $f_s$. That single sentence is
the sampling theorem; everything else is bookkeeping about whether the replicas collide.
\end{keybox}

\subsection{The Nyquist criterion and aliasing}
If $x$ is lowpass with bandwidth $W$ (so $X(f)=0$ for $|f|>W$), the replicas occupy
$[kf_s-W,\,kf_s+W]$. They do not overlap iff
\[
f_s>2W \qquad\text{(Nyquist rate $=2W$)}.
\]
Then $x(t)$ is recovered exactly by an ideal LPF of gain $T_s$ and cutoff between $W$ and $f_s-W$:
\[
x(t)=\sum_{n}x(nT_s)\,\sinc\!\left(\frac{t-nT_s}{T_s}\right)\qquad\text{(at $f_s=2W$).}
\]
If $f_s<2W$, replicas overlap and high frequencies fold down into the baseband as lower frequencies. This is
\textbf{aliasing}, and it is \emph{irreversible} --- no filter afterwards can undo it. Hence every practical
sampler is preceded by an \textbf{anti-aliasing filter}.

\begin{trapbox}
\textbf{The equality case.} At exactly $f_s=2W$ the replicas touch at $f=W$. This is fine for a strictly
bandlimited signal with no impulse at $\pm W$, but fails for e.g.\ $x(t)=\sin(2\pi Wt)$ sampled at $2W$: every
sample lands on a zero crossing and you recover nothing. Exam questions love this. Say ``$f_s>2W$, with
equality permissible only if $X(f)$ has no impulse at $\pm W$''.
\end{trapbox}

\subsection{Practical sampling}
\begin{itemize}[leftmargin=*]
\item \textbf{Natural sampling}: multiply by a pulse train of finite-width rectangles. Spectrum is
$\sum_k c_k X(f-kf_s)$ with $c_k$ the pulse train's Fourier coefficients --- replicas are weighted by a sinc
envelope but each is undistorted, so recovery is still exact.
\item \textbf{Flat-top sampling (PAM, sample-and-hold)}: each sample is held for width $\tau$. Now
$X_s(f)=\tau\sinc(f\tau)\cdot f_s\sum_kX(f-kf_s)$ --- the baseband copy is \emph{multiplied} by
$\sinc(f\tau)$, which droops across the band. This is \textbf{aperture distortion}, fixed by an equaliser of
response $1/\sinc(f\tau)$. This appears in EE341's PAM/PCM section.
\end{itemize}

\begin{worked}{8.1}
$x(t)=\cos(2\pi\cdot 300t)$ is sampled at $f_s=500$ Hz. What is heard after ideal lowpass reconstruction at
$250$ Hz?
\end{worked}
\soln
$W=300>f_s/2=250$, so aliasing. Replicas of the lines at $\pm300$ appear at $\pm300+k\cdot500$. Taking
$k=-1$: $300-500=-200$; taking $k=+1$: $-300+500=200$. So a pair of lines lands at $\pm200$ Hz inside the
reconstruction band. Output: a $200$ Hz tone. The alias frequency is $|f-kf_s|$ folded into $[0,f_s/2]$.

\section{Bandpass signals and the lowpass equivalent}

This is where your lecture notes currently are, and it is the conceptual heart of the first half of EE341.
Everything below is a direct consequence of Hermitian symmetry (\S5.4).

\subsection{The motivation}
A real bandpass signal at $f_c=900$ MHz with bandwidth $200$ kHz is a nightmare to simulate or analyse
directly --- you would need to sample at $\ge1.8$ GHz to represent a signal that carries only $200$ kHz of
information. The redundancy is twofold: (i) negative frequencies duplicate positive ones (Hermitian symmetry),
and (ii) the carrier itself carries no information. Strip both away and you get a \emph{lowpass} complex
signal of bandwidth $W$ that contains everything. That object is the complex envelope.

\subsection{Step 1: the analytic signal}
Delete negative frequencies:
\[
X_+(f)=X_{bp}(f)\,U(f),\qquad U(f)=\tfrac12\big(1+\sgn(f)\big).
\]
So $X_+(f)=\tfrac12 X_{bp}(f)+\tfrac12X_{bp}(f)\sgn(f)$. In the time domain, using
$\sgn(f)\FT^{-1}$-pairing with $\frac{1}{\ii\pi t}$, i.e.\ $\frac{1}{\pi t}\FT-\ii\sgn(f)$:
\[
x_+(t)=\frac12 x_{bp}(t)+\frac{\ii}{2}\,\hat x_{bp}(t),
\qquad
\hat x_{bp}(t):=x_{bp}(t)*\frac{1}{\pi t}
\]
where $\hat x$ is the \textbf{Hilbert transform} of $x$. (Your notes write
$x_+ = x_{bp}*\big(\tfrac12\delta(t)+\tfrac{1}{\ii2\pi t}\big)$ --- same thing; the
$\tfrac{1}{\ii 2\pi t}$ combines with the factor to give $\tfrac{\ii}{2}\hat x_{bp}$. Watch the $\ii$'s.)

\begin{keybox}
\textbf{The Hilbert transform is a $-90^\circ$ phase shifter.} Its frequency response is
$H(f)=-\ii\sgn(f)$: magnitude 1 everywhere, phase $-90^\circ$ for $f>0$ and $+90^\circ$ for $f<0$.
It does nothing to amplitudes; it only rotates phase. Key facts:
\[
\widehat{\cos(2\pi f_ct)}=\sin(2\pi f_ct),\qquad \widehat{\sin(2\pi f_ct)}=-\cos(2\pi f_ct),\qquad \hat{\hat x}=-x.
\]
\end{keybox}

\subsection{Step 2: shift down to baseband}
\[
x_\ell(t)=2x_+(t)e^{-\ii2\pi f_ct}
\qquad\text{(some texts, and your notes, omit the 2; be consistent)}
\]
Expanding with $x_+=\frac12(x_{bp}+\ii\hat x_{bp})$ and $e^{-\ii2\pi f_ct}=\cos-\ii\sin$:
\[
x_\ell(t)=x_I(t)+\ii\,x_Q(t)
\]
with
\begin{align}
x_I(t)&=x_{bp}(t)\cos(2\pi f_ct)+\hat x_{bp}(t)\sin(2\pi f_ct) &&\textbf{in-phase}\\
x_Q(t)&=\hat x_{bp}(t)\cos(2\pi f_ct)-x_{bp}(t)\sin(2\pi f_ct) &&\textbf{quadrature}
\end{align}
Both $x_I$ and $x_Q$ are \textbf{real lowpass signals of bandwidth $W$}. This is the block diagram on your
notes page 12.

\subsection{Step 3: reconstruct}
\begin{keybox}
\[
\boxed{\;x_{bp}(t)=\operatorname{Re}\big\{x_\ell(t)\,e^{\ii2\pi f_ct}\big\}
= x_I(t)\cos(2\pi f_ct)-x_Q(t)\sin(2\pi f_ct)\;}
\]
\end{keybox}
That minus sign is not optional --- it comes from $\operatorname{Re}\{(x_I+\ii x_Q)(\cos+\ii\sin)\}$. Losing
it is the most common error in this topic.

\subsection{Envelope and phase}
Write $x_\ell(t)=A(t)e^{\ii\phi(t)}$ with
\[
A(t)=\sqrt{x_I^2(t)+x_Q^2(t)}\ \ (\ge0),
\qquad
\phi(t)=\tan^{-1}\!\frac{x_Q(t)}{x_I(t)} .
\]
Then $x_{bp}(t)=A(t)\cos\big(2\pi f_ct+\phi(t)\big)$. $A(t)$ is the \textbf{envelope} --- and it is exactly
what an envelope detector (diode + RC) outputs.

\begin{keybox}
\textbf{Now the AM results become one-liners.}
\begin{itemize}[nosep,leftmargin=*]
\item \textbf{DSB-SC:} $x_c(t)=Ax(t)\cos2\pi f_ct=\operatorname{Re}\{Ax(t)e^{\ii2\pi f_ct}\}$, so
$x_\ell=Ax(t)$, which is real. Envelope $=A|x(t)|$. Since $x(t)$ goes negative, the envelope is the
\emph{rectified} message --- so an envelope detector \textbf{fails}, and the sign flips show up as the
$180^\circ$ phase reversals drawn on your notes page 3.
\item \textbf{DSB-AM:} $x_c'(t)=A(1+\mu x(t))\cos2\pi f_ct$, so $x_\ell=A(1+\mu x(t))$. If $\mu|x(t)|<1$ for
all $t$ (i.e.\ modulation index under 100\%), then $x_\ell>0$ always, the envelope is
$A(1+\mu x(t))$, and an envelope detector recovers $x(t)$ up to a DC offset. \textbf{Over-modulation}
($\mu|x|>1$) makes $x_\ell$ change sign and the envelope detector produces the rectified, distorted mess.
\item \textbf{Angle modulation:} $A(t)=$ const, all the information is in $\phi(t)$.
\end{itemize}
\end{keybox}

\subsection{Noise in the same language}
Your notes page 11 does this: pass bandpass noise $n(t)$ through the receiver's BPF and write
$n(t)=n_I(t)\cos2\pi f_ct-n_Q(t)\sin2\pi f_ct$. Then received signal + noise has envelope
\[
A(t)=\sqrt{\big(x_I(t)+n_I(t)\big)^2+\big(x_Q(t)+n_Q(t)\big)^2},
\]
which is the starting point for every SNR calculation in analog communication. For white Gaussian noise
filtered to bandwidth $2W$, $n_I$ and $n_Q$ turn out to be independent, zero-mean, each with the \emph{same}
variance as $n$ --- proved with the tools of \S10.

\begin{worked}{9.1}
Find the lowpass equivalent of $x_{bp}(t)=\cos\big(2\pi(f_c+\Delta f)t\big)$ with respect to carrier $f_c$.
\end{worked}
\soln
$x_{bp}=\operatorname{Re}\{e^{\ii2\pi(f_c+\Delta f)t}\}=\operatorname{Re}\{e^{\ii2\pi\Delta f t}e^{\ii2\pi f_ct}\}$,
so $x_\ell(t)=e^{\ii2\pi\Delta ft}$, giving $x_I=\cos(2\pi\Delta ft)$, $x_Q=\sin(2\pi\Delta ft)$, envelope
$A(t)=1$. \textbf{A frequency offset becomes a slow rotation of the complex envelope} --- which is precisely
the picture on your notes page 11 where $x_\ell(t)=Ax(t)+Ae^{\ii2\pi\Delta ft}$ and hence
$x_I=Ax(t)+A\cos(2\pi\Delta ft)$, $x_Q=A\sin(2\pi\Delta ft)$. Everything about carrier offset, phase noise,
and synchronisation is easier to see in this representation.

\begin{worked}{9.2}
Show the bandpass filter $h_{bp}(t)=2\operatorname{Re}\{h_\ell(t)e^{\ii2\pi f_ct}\}$ acting on
$x_{bp}$ produces output whose complex envelope is $y_\ell=x_\ell * h_\ell$.
\end{worked}
\soln
Sketch: in the frequency domain, $Y_{bp}(f)=H_{bp}(f)X_{bp}(f)$. Restrict to $f>0$, where
$X_{bp}(f)=\tfrac12X_\ell(f-f_c)$ and $H_{bp}(f)=H_\ell(f-f_c)$. Multiply, shift back by $f_c$, and read off
$Y_\ell(f)=X_\ell(f)H_\ell(f)$; invert. \textbf{Why this matters:} it means you can simulate an entire
bandpass communication link --- modulator, channel filter, receiver filter, noise --- entirely at baseband
with complex signals, at a sampling rate set by $W$ and not by $f_c$. That is how every real simulator (and
every software-defined radio) works.
%==============================================================
\section*{\color{keyblue}PART D --- Random Processes Primer}
\addcontentsline{toc}{section}{PART D --- Random Processes Primer}
%==============================================================

\section{Just enough randomness to survive the noise chapter}

Skim this before the quiz; return to it properly before the mid-semester. The syllabus lists it as a prereq
because noise is not a signal you can write down --- you can only describe its statistics.

\subsection{Random variables, in one box}
\begin{keybox}
Mean $\mu_X=E[X]$; variance $\sigma_X^2=E[(X-\mu_X)^2]=E[X^2]-\mu_X^2$.\\
Gaussian: $p_X(x)=\dfrac{1}{\sqrt{2\pi\sigma^2}}e^{-(x-\mu)^2/2\sigma^2}$.\\
The Q-function: $Q(x)=P(X>x)$ for standard Gaussian $=\displaystyle\int_x^\infty\frac{1}{\sqrt{2\pi}}e^{-t^2/2}\dd t$.
\end{keybox}
$Q(\cdot)$ is the function every digital-communication error probability is expressed in. Get comfortable with
$Q(x)=1-Q(-x)$ and its rapid decay. You will meet $P_e=Q\!\left(\sqrt{2E_b/N_0}\right)$ for BPSK.

\subsection{Random processes}
A random process $X(t)$ is a family of random variables, one per time instant. Two statistics matter:
\[
\mu_X(t)=E[X(t)],\qquad R_X(t_1,t_2)=E[X(t_1)X(t_2)]\quad\text{(autocorrelation)}.
\]

\textbf{Wide-sense stationary (WSS)} means: $\mu_X(t)=\mu_X$ is constant, and $R_X(t_1,t_2)$ depends only on
the lag $\tau=t_1-t_2$, so we write $R_X(\tau)$. Almost every noise model in this course is WSS.

Properties of $R_X(\tau)$ for a WSS process:
\begin{itemize}[nosep,leftmargin=*]
\item $R_X(0)=E[X^2(t)]=$ average power.
\item $R_X(-\tau)=R_X(\tau)$ (even).
\item $|R_X(\tau)|\le R_X(0)$.
\end{itemize}

\subsection{Power spectral density and the Wiener--Khinchin theorem}
\begin{keybox}
\[
\boxed{S_X(f)=\int_{-\infty}^{\infty}R_X(\tau)e^{-\ii2\pi f\tau}\dd\tau}
\qquad\text{and}\qquad
P_X=R_X(0)=\int_{-\infty}^{\infty}S_X(f)\dd f .
\]
\textbf{The PSD is the Fourier transform of the autocorrelation.} This is what replaces $|X(f)|^2$ when the
signal is a power signal with no ordinary Fourier transform.
\end{keybox}
$S_X(f)$ is real, even, and non-negative. It has units W/Hz.

\subsection{The one result you will use constantly}
\begin{keybox}
\textbf{WSS process through an LTI filter.} If $Y=X*h$, then
\[
S_Y(f)=|H(f)|^2 S_X(f),\qquad \mu_Y=\mu_X H(0).
\]
\end{keybox}
Note $|H(f)|^2$, \emph{not} $H(f)$: phase information is destroyed, because the PSD does not know about phase.

\subsection{White noise}
\textbf{AWGN} (additive white Gaussian noise) has
\[
S_N(f)=\frac{N_0}{2}\ \text{(W/Hz, two-sided, for all $f$)},\qquad R_N(\tau)=\frac{N_0}{2}\delta(\tau).
\]
Two consequences: (i) samples at distinct times are uncorrelated (and, being Gaussian, independent);
(ii) it has infinite power, so it is a fiction --- but a harmless one, because every receiver filters it.
After an ideal BPF of bandwidth $2W$ (both sidebands, at $\pm f_c$), the noise power is
\[
P_N=\int|H(f)|^2\frac{N_0}{2}\dd f=\frac{N_0}{2}\cdot 2\cdot(2W)=2N_0W .
\]
Combine this with the signal power from Worked Problem 1.1 and you have the SNR expressions your notes are
building towards.

\begin{worked}{10.1}
Show that for bandpass noise $n(t)=n_I(t)\cos2\pi f_ct-n_Q(t)\sin2\pi f_ct$ with $S_N(f)$ flat of height
$N_0/2$ over $|f\mp f_c|<W$, the components satisfy $\sigma_{n_I}^2=\sigma_{n_Q}^2=\sigma_n^2$.
\end{worked}
\soln
Sketch of the standard argument: $n_I$ and $n_Q$ are obtained from $n$ by the operations in \S9.3, which are
lowpass filters of one-sided width $W$ applied to frequency-shifted copies. Their PSDs come out as
$S_{n_I}(f)=S_{n_Q}(f)=S_N(f-f_c)+S_N(f+f_c)$ for $|f|<W$, i.e.\ flat of height $N_0$ over $[-W,W]$. So
\[
\sigma_{n_I}^2=\sigma_{n_Q}^2=N_0\cdot 2W = 2N_0W=\sigma_n^2 .
\]
Equal variance to the bandpass noise itself --- this is the fact that makes envelope-detector and
coherent-detector SNR comparisons work out.

%==============================================================
\section{Problem set (do this last, timed)}
%==============================================================

Give yourself \textbf{60 minutes}. Solutions follow --- do not look until you have finished.

\begin{enumerate}[leftmargin=*,itemsep=6pt]
\item Sketch $|X(f)|$ for $x(t)=\rect(t/T)\cos(2\pi f_ct)$ with $f_c\gg1/T$. What is the bandwidth?
\item A real signal $x(t)$ has $X(f)$ with $X(f)=0$ for $|f|>5$ kHz. It is multiplied by $\cos(2\pi\cdot100{,}000\,t)$.
Over what frequency range does the result live? What is the minimum sampling rate needed to represent the
\emph{result} without loss if you sample it directly as a bandpass signal? (Hint: think about the lowpass
equivalent.)
\item Prove that if $x(t)$ is real and even, $X(f)$ is real and even.
\item Compute $\big(e^{-at}u(t)\big)*\big(e^{-bt}u(t)\big)$ for $a\neq b$, two ways: directly, and via
Fourier transforms with partial fractions.
\item A DSB-AM signal has $A=1$, $\mu=0.5$, and message $x(t)=\cos(2\pi\cdot1000t)$. Compute the total power,
the sideband power, and the power efficiency.
\item Find the Hilbert transform of $x(t)=\cos(2\pi f_0t)+\sin(2\pi f_0 t)$, and hence write its analytic signal.
\item An LTI system has $h(t)=\delta(t)-\delta(t-T)$. Find $H(f)$, sketch $|H(f)|$, and state at which
frequencies the system nulls the input. (This is a two-ray multipath channel --- it shows up in the ISI
chapter.)
\item $x(t)$ is bandlimited to $W=4$ kHz and sampled at $f_s=7$ kHz. If $x$ contains a tone at $3.5$ kHz, at
what frequency does its alias appear after reconstruction with an ideal $3.5$ kHz LPF?
\end{enumerate}

\subsection*{Solutions}

\textbf{1.} By the modulation property, $X(f)=\frac{T}{2}\sinc\big((f-f_c)T\big)+\frac{T}{2}\sinc\big((f+f_c)T\big)$:
two sinc lobes centred at $\pm f_c$. First nulls at $f_c\pm 1/T$, so the null-to-null bandwidth is $2/T$.
(The baseband pulse had bandwidth $1/T$; modulation doubles it. Every DSB scheme doubles bandwidth.)

\medskip
\textbf{2.} The product occupies $95$--$105$ kHz and $-105$ to $-95$ kHz: a bandpass signal with $f_c=100$ kHz,
$W=5$ kHz, total bandwidth $10$ kHz. Naive lowpass reasoning would demand $f_s>2\times105=210$ kHz. But the
lowpass equivalent $x_\ell$ is complex with bandwidth $5$ kHz, so $10$ kHz complex samples/second suffice ---
i.e.\ $10$ kHz on each of the I and Q branches, $20$ k real samples/second. This is \emph{bandpass sampling},
and it is the practical reason \S9 exists.

\medskip
\textbf{3.} $X(f)=\int x(t)e^{-\ii2\pi ft}\dd t=\int x(t)\cos(2\pi ft)\dd t-\ii\int x(t)\sin(2\pi ft)\dd t$.
The second integrand is (even)$\times$(odd) $=$ odd, so it integrates to zero over $\mathbb{R}$. Hence $X(f)$
is real. And $X(-f)=\int x(t)\cos(2\pi ft)\dd t=X(f)$ since cosine is even, so $X$ is even. $\square$

\medskip
\textbf{4.} Direct: for $t>0$,
$\int_0^t e^{-a\tau}e^{-b(t-\tau)}\dd\tau=e^{-bt}\int_0^te^{(b-a)\tau}\dd\tau=\dfrac{e^{-at}-e^{-bt}}{b-a}$,
and $0$ for $t<0$.\\
Via FT: the product is $\dfrac{1}{(a+\ii2\pi f)(b+\ii2\pi f)}=\dfrac{1}{b-a}\left[\dfrac{1}{a+\ii2\pi f}-\dfrac{1}{b+\ii2\pi f}\right]$,
which inverts term by term to the same answer. The second route is the one to use under time pressure.

\medskip
\textbf{5.} $P_x=\frac12$ for a unit-amplitude cosine. Total power
$P=\frac{A^2}{2}(1+\mu^2P_x)=\frac12\left(1+0.25\cdot0.5\right)=\frac12(1.125)=0.5625$ W.
Carrier power $=\frac{A^2}{2}=0.5$ W; sideband (message) power $=\frac{A^2\mu^2}{2}P_x=0.0625$ W.
Efficiency $\eta=0.0625/0.5625=1/9\approx11.1\%$. Even at $\mu=1$ with a sinusoidal message the ceiling is
$33.3\%$ --- two-thirds of your transmitter power goes into a carrier that carries no information. That is the
argument for DSB-SC.

\medskip
\textbf{6.} $\hat x(t)=\sin(2\pi f_0t)-\cos(2\pi f_0t)$ (using $\hat{\cos}=\sin$, $\hat{\sin}=-\cos$).
Then $x_+(t)=x+\ii\hat x=(1-\ii)\cos(2\pi f_0t)+(1+\ii)\sin(2\pi f_0t)$. Alternatively note
$x(t)=\sqrt2\cos(2\pi f_0t-\pi/4)$, so $x_+(t)=\sqrt2\,e^{\ii(2\pi f_0t-\pi/4)}$ --- much cleaner, and it
confirms the analytic signal of a sinusoid is just the corresponding rotating phasor.
(Convention note: with the $x_+=\tfrac12(x+\ii\hat x)$ definition used in \S9, halve these.)

\medskip
\textbf{7.} $H(f)=1-e^{-\ii2\pi fT}$. Factor: $=e^{-\ii\pi fT}\big(e^{\ii\pi fT}-e^{-\ii\pi fT}\big)=2\ii\sin(\pi fT)e^{-\ii\pi fT}$,
so $|H(f)|=2|\sin(\pi fT)|$. Nulls where $\sin(\pi fT)=0$, i.e.\ $f=k/T$ for integer $k$ (including DC).
A comb of nulls spaced $1/T$ apart --- these are the ``frequency-selective fades'' of a two-path channel.

\medskip
\textbf{8.} $f_s=7$ kHz means replicas at $\pm3.5+7k$. The tone at $3.5$ kHz sits exactly at $f_s/2$: the
replica from $k=-1$ lands at $3.5-7=-3.5$ kHz, coinciding with the original negative-frequency line. This is
the boundary case of the trap in \S8.2 --- the tone is not aliased to a different frequency, but its amplitude
and phase are not recoverable in general (samples may all land on zero crossings). Note also that
$f_s=7<2W=8$ kHz, so the signal is undersampled overall and other components between $3.5$ and $4$ kHz
\emph{are} genuinely aliased into the band.

%==============================================================
\section{One-page recall sheet}
%==============================================================

\begin{keybox}
\textbf{Transforms}\quad $X(f)=\int x e^{-\ii2\pi ft}\dd t$; \ $x(t)=\int Xe^{\ii2\pi ft}\dd f$

\textbf{Convolution} $x*h\FT XH$ \quad\textbf{Product} $xh\FT X*H$

\textbf{Shift} $x(t-t_0)\FT Xe^{-\ii2\pi ft_0}$ \quad $xe^{\ii2\pi f_ct}\FT X(f-f_c)$

\textbf{Modulation} $x\cos2\pi f_ct\FT \tfrac12X(f-f_c)+\tfrac12X(f+f_c)$

\textbf{Scaling} $x(at)\FT\frac{1}{|a|}X(f/a)$ \quad\textbf{Duality} $X(t)\FT x(-f)$

\textbf{Real $x$} $\Rightarrow X(-f)=X^*(f)$

\textbf{Parseval} $\int|x|^2\dd t=\int|X|^2\dd f$ \quad $X(0)=\int x\dd t$, \ $x(0)=\int X\dd f$

\medskip
\textbf{Pairs} $\rect(t/T)\FT T\sinc(fT)$; \ $\tri(t/T)\FT T\sinc^2(fT)$; \
$\cos2\pi f_ct\FT\tfrac12[\delta(f-f_c)+\delta(f+f_c)]$; \
$\frac{1}{\pi t}\FT-\ii\sgn f$; \ $\sum_n\delta(t-nT_s)\FT f_s\sum_k\delta(f-kf_s)$

\medskip
\textbf{Sampling} $X_s(f)=f_s\sum_kX(f-kf_s)$; \ need $f_s>2W$

\textbf{Bandpass} $x_{bp}=\operatorname{Re}\{x_\ell e^{\ii2\pi f_ct}\}=x_I\cos2\pi f_ct-x_Q\sin2\pi f_ct$;
\ $A=\sqrt{x_I^2+x_Q^2}$

\textbf{Coherent detection} output $\propto\cos\phi$ --- zero at $\phi=90^\circ$

\textbf{DSB-SC} $x_\ell=Ax(t)$ (real, sign-changing $\Rightarrow$ envelope detection fails)

\textbf{DSB-AM} $x_\ell=A(1+\mu x)$; envelope detection works iff $\mu|x|<1$; $\eta=\frac{\mu^2P_x}{1+\mu^2P_x}$

\textbf{Noise} $S_Y=|H|^2S_X$; \ white: $S_N=N_0/2$, $R_N=\frac{N_0}{2}\delta(\tau)$
\end{keybox}

\vspace{1em}
\begin{trapbox}
\textbf{The five mistakes that cost the most marks}
\begin{enumerate}[nosep,leftmargin=*]
\item Mixing $f$ and $\omega$ conventions mid-problem (stray $2\pi$'s).
\item Forgetting the minus sign in $x_{bp}=x_I\cos-x_Q\sin$.
\item Calling the bandwidth of a real lowpass signal $2W$ when the course means $W$ (or vice versa) --- state your convention.
\item Applying the modulation property without checking $f_c>W$, so the shifted copies actually don't overlap.
\item Using $|X(f)|^2$ (energy spectral density) for a modulated carrier or noise, which is a \emph{power} signal and needs PSD.
\end{enumerate}
\end{trapbox}

\vfill
\noindent\small\textit{Built against the EE341 (Autumn 2026) course introduction and the handwritten lecture
notes in your course folder, so the notation matches what Prof.\ Appaiah uses. Reference texts if you want
more depth: Madhow, \emph{Introduction to Communication Systems}, Ch.~2 (this is the best match for the
lowpass-equivalent material in \S9); Lathi, Ch.~3--4 for signals and Fourier; Haykin \& Moher Ch.~2 for the
bandpass/Hilbert treatment.}

\end{document}
